% assets/w06-alloc-geometry.svg — 추진기 셋이 앉은 자리와 모멘트 팔
%
%   bash _tools/tikz2svg.sh assets/src/w06-alloc-geometry.tex assets/w06-alloc-geometry.svg
%
% 이 그림이 말하는 것
%   배분행렬 T 의 숫자는 외우는 것이 아니라 **추진기가 어디 앉아 있는지를 적은 것**이다.
%   한 추진기가 힘 F 를 내면 그 힘은
%       앞뒤로 미는 몫 X,  옆으로 미는 몫 Y,  배를 돌리는 몫 N = (모멘트 팔) x F
%   로 갈린다. 팔의 길이가 곧 행렬의 세 번째 줄이다.
%
%   좌우 추진기  팔 = b  (중심선에서 좌우로 떨어진 거리, half_beam = 1.027135 m)
%   선수 추진기  팔 = l3 (중심에서 앞으로 떨어진 거리, l_bow = 0.75 m)
%
% 좌표 — 배를 위에서 본다. 선체 좌표계
%   x_b = 선수 방향(그림 위), y_b = 우현 방향(그림 오른쪽)
%   \BP{xb}{yb} 로 적는다. 1 m = 20 mm
%   화면 x = y_b, 화면 y = x_b 이므로 매크로 안에서 인자 순서가 뒤집힌다
%
% 겹침을 피하려고 정한 것
%   돌림(호)은 선체 밖 반지름 48 mm 에 그린다 — 선체 반폭이 20.5 mm 라 닿지 않는다
%   양의 요 N 은 x_b 가 y_b 쪽으로 도는 것 = 화면에서 시계방향 = arc 각이 **줄어드는** 쪽

\documentclass[border=6pt]{standalone}
% capstone-style.tex — 이 볼트의 TikZ 그림이 공유하는 스타일
%
% 쓰는 법:  % capstone-style.tex — 이 볼트의 TikZ 그림이 공유하는 스타일
%
% 쓰는 법:  % capstone-style.tex — 이 볼트의 TikZ 그림이 공유하는 스타일
%
% 쓰는 법:  \input{capstone-style}   (assets/src/ 안에서)
% 빌드:     bash _tools/tikz2svg.sh assets/src/w02-node-topic.tex assets/w02-node-topic.svg
%
% 왜 TikZ 인가
%   손으로 SVG 좌표를 쓰면 화살촉·선 굵기·글자 크기가 그림마다 미묘하게 달라진다.
%   그것이 "자료가 어설퍼 보인다"의 정체다. 스타일을 한 번만 정하고 상대좌표로
%   배치하면 좌표를 고쳐도 그림이 무너지지 않는다.
%   수식은 본문과 같은 TeX 글꼴로 나온다.
%
% 한글
%   ko.TeX(DVI 경로)를 쓴다. XeLaTeX 은 TikZ 그래픽을 PDF special 로 내보내는데
%   dvisvgm 이 그것을 받으려면 Ghostscript < 10.01 이 필요하다. 이 PC 에는
%   10.07 이 깔려 있어 그 경로가 막힌다. latex -> DVI -> dvisvgm 이 유일하게
%   동작하는 조합이다.

\usepackage[hangul]{kotex}
\usepackage{tikz}
\usepackage{amsmath}
\usepackage{xcolor}
\usetikzlibrary{arrows.meta, positioning, calc, fit, backgrounds, decorations.pathmorphing}

% ---- 색 --------------------------------------------------------------------
% 강조는 하나만. 나머지는 검정. 흑백 인쇄에서도 읽혀야 한다.
\definecolor{ink}{HTML}{1A1A1A}
\definecolor{soft}{HTML}{666666}
\definecolor{accent}{HTML}{7C3AED}
\definecolor{warn}{HTML}{B91C1C}

% 계층 색 — Simulink 모델의 블록 색과 같은 규칙을 쓴다
\definecolor{guid}{HTML}{1D4ED8}    % 유도
\definecolor{ctrl}{HTML}{EA580C}    % 제어
\definecolor{thr}{HTML}{D97706}     % 추진기
\definecolor{plant}{HTML}{16A34A}   % 운동모델
\definecolor{nav}{HTML}{7C3AED}     % 항법 · 신호처리
\definecolor{pathc}{HTML}{0D9488}   % 계획 경로 · 항적
\definecolor{errc}{HTML}{C2410C}    % 오차
\definecolor{flow}{HTML}{0369A1}    % 조류 · 바람

% ---- 선과 화살촉 -----------------------------------------------------------
\tikzset{
  sig/.style   = {ink, line width=0.5pt, -{Stealth[length=4pt, width=3pt]}},
  wire/.style  = {ink, line width=0.5pt},
  fb/.style    = {ink, line width=0.5pt, densely dashed,
                  -{Stealth[length=4pt, width=3pt]}},
  asig/.style  = {accent, line width=0.5pt, densely dashed,
                  -{Stealth[length=4pt, width=3pt]}},
  % 블록 — 흰 바탕, 직각, 검은 테두리. 색은 테두리에만 준다
  blk/.style   = {draw=ink, line width=0.55pt, fill=white,
                  minimum width=18mm, minimum height=8.5mm, inner sep=2pt, font=\small},
  gblk/.style  = {blk, draw=guid},
  cblk/.style  = {blk, draw=ctrl},
  tblk/.style  = {blk, draw=thr},
  pblk/.style  = {blk, draw=plant},
  nblk/.style  = {blk, draw=nav},
  % 합산점 · 분기점
  sum/.style   = {draw=ink, line width=0.55pt, fill=white, circle,
                  inner sep=0pt, minimum size=4.6mm, font=\scriptsize},
  dot/.style   = {ink, fill=ink, circle, inner sep=0pt, minimum size=1.5mm},
  % 글자
  sname/.style = {font=\small, inner sep=1.5pt},
  note/.style  = {font=\scriptsize, soft, inner sep=1.5pt},
  wnote/.style = {font=\scriptsize, warn, inner sep=1.5pt},
  anote/.style = {font=\scriptsize, accent, inner sep=1.5pt},
  axis/.style  = {ink, line width=0.5pt},
  tick/.style  = {font=\scriptsize, soft},
  % 묶음 상자 — 서브시스템 표시
  grp/.style   = {draw=soft, line width=0.4pt, densely dotted, rounded corners=1mm,
                  inner sep=3mm},
}
   (assets/src/ 안에서)
% 빌드:     bash _tools/tikz2svg.sh assets/src/w02-node-topic.tex assets/w02-node-topic.svg
%
% 왜 TikZ 인가
%   손으로 SVG 좌표를 쓰면 화살촉·선 굵기·글자 크기가 그림마다 미묘하게 달라진다.
%   그것이 "자료가 어설퍼 보인다"의 정체다. 스타일을 한 번만 정하고 상대좌표로
%   배치하면 좌표를 고쳐도 그림이 무너지지 않는다.
%   수식은 본문과 같은 TeX 글꼴로 나온다.
%
% 한글
%   ko.TeX(DVI 경로)를 쓴다. XeLaTeX 은 TikZ 그래픽을 PDF special 로 내보내는데
%   dvisvgm 이 그것을 받으려면 Ghostscript < 10.01 이 필요하다. 이 PC 에는
%   10.07 이 깔려 있어 그 경로가 막힌다. latex -> DVI -> dvisvgm 이 유일하게
%   동작하는 조합이다.

\usepackage[hangul]{kotex}
\usepackage{tikz}
\usepackage{amsmath}
\usepackage{xcolor}
\usetikzlibrary{arrows.meta, positioning, calc, fit, backgrounds, decorations.pathmorphing}

% ---- 색 --------------------------------------------------------------------
% 강조는 하나만. 나머지는 검정. 흑백 인쇄에서도 읽혀야 한다.
\definecolor{ink}{HTML}{1A1A1A}
\definecolor{soft}{HTML}{666666}
\definecolor{accent}{HTML}{7C3AED}
\definecolor{warn}{HTML}{B91C1C}

% 계층 색 — Simulink 모델의 블록 색과 같은 규칙을 쓴다
\definecolor{guid}{HTML}{1D4ED8}    % 유도
\definecolor{ctrl}{HTML}{EA580C}    % 제어
\definecolor{thr}{HTML}{D97706}     % 추진기
\definecolor{plant}{HTML}{16A34A}   % 운동모델
\definecolor{nav}{HTML}{7C3AED}     % 항법 · 신호처리
\definecolor{pathc}{HTML}{0D9488}   % 계획 경로 · 항적
\definecolor{errc}{HTML}{C2410C}    % 오차
\definecolor{flow}{HTML}{0369A1}    % 조류 · 바람

% ---- 선과 화살촉 -----------------------------------------------------------
\tikzset{
  sig/.style   = {ink, line width=0.5pt, -{Stealth[length=4pt, width=3pt]}},
  wire/.style  = {ink, line width=0.5pt},
  fb/.style    = {ink, line width=0.5pt, densely dashed,
                  -{Stealth[length=4pt, width=3pt]}},
  asig/.style  = {accent, line width=0.5pt, densely dashed,
                  -{Stealth[length=4pt, width=3pt]}},
  % 블록 — 흰 바탕, 직각, 검은 테두리. 색은 테두리에만 준다
  blk/.style   = {draw=ink, line width=0.55pt, fill=white,
                  minimum width=18mm, minimum height=8.5mm, inner sep=2pt, font=\small},
  gblk/.style  = {blk, draw=guid},
  cblk/.style  = {blk, draw=ctrl},
  tblk/.style  = {blk, draw=thr},
  pblk/.style  = {blk, draw=plant},
  nblk/.style  = {blk, draw=nav},
  % 합산점 · 분기점
  sum/.style   = {draw=ink, line width=0.55pt, fill=white, circle,
                  inner sep=0pt, minimum size=4.6mm, font=\scriptsize},
  dot/.style   = {ink, fill=ink, circle, inner sep=0pt, minimum size=1.5mm},
  % 글자
  sname/.style = {font=\small, inner sep=1.5pt},
  note/.style  = {font=\scriptsize, soft, inner sep=1.5pt},
  wnote/.style = {font=\scriptsize, warn, inner sep=1.5pt},
  anote/.style = {font=\scriptsize, accent, inner sep=1.5pt},
  axis/.style  = {ink, line width=0.5pt},
  tick/.style  = {font=\scriptsize, soft},
  % 묶음 상자 — 서브시스템 표시
  grp/.style   = {draw=soft, line width=0.4pt, densely dotted, rounded corners=1mm,
                  inner sep=3mm},
}
   (assets/src/ 안에서)
% 빌드:     bash _tools/tikz2svg.sh assets/src/w02-node-topic.tex assets/w02-node-topic.svg
%
% 왜 TikZ 인가
%   손으로 SVG 좌표를 쓰면 화살촉·선 굵기·글자 크기가 그림마다 미묘하게 달라진다.
%   그것이 "자료가 어설퍼 보인다"의 정체다. 스타일을 한 번만 정하고 상대좌표로
%   배치하면 좌표를 고쳐도 그림이 무너지지 않는다.
%   수식은 본문과 같은 TeX 글꼴로 나온다.
%
% 한글
%   ko.TeX(DVI 경로)를 쓴다. XeLaTeX 은 TikZ 그래픽을 PDF special 로 내보내는데
%   dvisvgm 이 그것을 받으려면 Ghostscript < 10.01 이 필요하다. 이 PC 에는
%   10.07 이 깔려 있어 그 경로가 막힌다. latex -> DVI -> dvisvgm 이 유일하게
%   동작하는 조합이다.

\usepackage[hangul]{kotex}
\usepackage{tikz}
\usepackage{amsmath}
\usepackage{xcolor}
\usetikzlibrary{arrows.meta, positioning, calc, fit, backgrounds, decorations.pathmorphing}

% ---- 색 --------------------------------------------------------------------
% 강조는 하나만. 나머지는 검정. 흑백 인쇄에서도 읽혀야 한다.
\definecolor{ink}{HTML}{1A1A1A}
\definecolor{soft}{HTML}{666666}
\definecolor{accent}{HTML}{7C3AED}
\definecolor{warn}{HTML}{B91C1C}

% 계층 색 — Simulink 모델의 블록 색과 같은 규칙을 쓴다
\definecolor{guid}{HTML}{1D4ED8}    % 유도
\definecolor{ctrl}{HTML}{EA580C}    % 제어
\definecolor{thr}{HTML}{D97706}     % 추진기
\definecolor{plant}{HTML}{16A34A}   % 운동모델
\definecolor{nav}{HTML}{7C3AED}     % 항법 · 신호처리
\definecolor{pathc}{HTML}{0D9488}   % 계획 경로 · 항적
\definecolor{errc}{HTML}{C2410C}    % 오차
\definecolor{flow}{HTML}{0369A1}    % 조류 · 바람

% ---- 선과 화살촉 -----------------------------------------------------------
\tikzset{
  sig/.style   = {ink, line width=0.5pt, -{Stealth[length=4pt, width=3pt]}},
  wire/.style  = {ink, line width=0.5pt},
  fb/.style    = {ink, line width=0.5pt, densely dashed,
                  -{Stealth[length=4pt, width=3pt]}},
  asig/.style  = {accent, line width=0.5pt, densely dashed,
                  -{Stealth[length=4pt, width=3pt]}},
  % 블록 — 흰 바탕, 직각, 검은 테두리. 색은 테두리에만 준다
  blk/.style   = {draw=ink, line width=0.55pt, fill=white,
                  minimum width=18mm, minimum height=8.5mm, inner sep=2pt, font=\small},
  gblk/.style  = {blk, draw=guid},
  cblk/.style  = {blk, draw=ctrl},
  tblk/.style  = {blk, draw=thr},
  pblk/.style  = {blk, draw=plant},
  nblk/.style  = {blk, draw=nav},
  % 합산점 · 분기점
  sum/.style   = {draw=ink, line width=0.55pt, fill=white, circle,
                  inner sep=0pt, minimum size=4.6mm, font=\scriptsize},
  dot/.style   = {ink, fill=ink, circle, inner sep=0pt, minimum size=1.5mm},
  % 글자
  sname/.style = {font=\small, inner sep=1.5pt},
  note/.style  = {font=\scriptsize, soft, inner sep=1.5pt},
  wnote/.style = {font=\scriptsize, warn, inner sep=1.5pt},
  anote/.style = {font=\scriptsize, accent, inner sep=1.5pt},
  axis/.style  = {ink, line width=0.5pt},
  tick/.style  = {font=\scriptsize, soft},
  % 묶음 상자 — 서브시스템 표시
  grp/.style   = {draw=soft, line width=0.4pt, densely dotted, rounded corners=1mm,
                  inner sep=3mm},
}


\begin{document}
\begin{tikzpicture}[x=1mm, y=1mm]

\pgfmathsetmacro{\BS}{20}                        % 1 m = 20 mm
\newcommand{\BP}[2]{({(#2)*\BS},{(#1)*\BS})}     % #1 = x_b (선수), #2 = y_b (우현)

\pgfmathsetmacro{\bb}{1.027135}                  % half_beam [m]
\pgfmathsetmacro{\lb}{0.75}                      % l_bow [m]

% --- 선체 외곽선 + 추진기 자리 ---------------------------------------------
\newcommand{\hullplan}{%
  \draw[ink, line width=0.6pt, fill=white]
       \BP{-2.0}{-\bb} -- \BP{1.3}{-\bb} -- \BP{2.2}{0} -- \BP{1.3}{\bb}
       -- \BP{-2.0}{\bb} -- cycle;
  \draw[soft, line width=0.3pt, densely dashed] \BP{-2.1}{0} -- \BP{2.3}{0};
  % 선체 좌표축
  \draw[nav, line width=0.6pt, -{Stealth[length=3.4pt,width=2.4pt]}]
       \BP{0}{0} -- \BP{1.95}{0};
  \node[font=\scriptsize, nav, anchor=south east] at \BP{1.98}{-0.1} {$x_b$ (선수)};
  \draw[nav, line width=0.6pt, -{Stealth[length=3.4pt,width=2.4pt]}]
       \BP{0}{0} -- \BP{0}{0.82};
  \node[font=\scriptsize, nav, anchor=north east] at \BP{-0.06}{0.88} {$y_b$ (우현)};
  \node[dot] at \BP{0}{0} {};
  \node[font=\scriptsize, soft, anchor=north east] at \BP{-0.1}{-0.1} {무게중심};
  % 추진기 자리
  \node[circle, draw=thr, fill=white, line width=0.7pt, minimum size=3.4mm,
        inner sep=0pt] at \BP{-1.8}{-\bb} {};
  \node[circle, draw=thr, fill=white, line width=0.7pt, minimum size=3.4mm,
        inner sep=0pt] at \BP{-1.8}{\bb} {};
  \node[circle, draw=thr, fill=white, line width=0.7pt, minimum size=3.4mm,
        inner sep=0pt] at \BP{\lb}{0} {};}

% ===========================================================================
% (가) 추진기가 앉은 자리와 팔
% ===========================================================================
\begin{scope}[shift={(0,0)}]
  \node[font=\small\bfseries, anchor=north west] at (-50,66)
       {(가) 세 추진기가 앉은 자리 — 배분행렬의 숫자는 여기서 나온다};

  \hullplan

  % 힘 화살표
  \draw[thr, line width=1.2pt, -{Stealth[length=4pt,width=3pt]}]
       \BP{-1.8}{-\bb} -- \BP{-0.6}{-\bb};
  \node[font=\scriptsize, thr, anchor=north east] at \BP{-1.95}{-\bb+0.05} {$F_L$};
  \draw[thr, line width=1.2pt, -{Stealth[length=4pt,width=3pt]}]
       \BP{-1.8}{\bb} -- \BP{-0.6}{\bb};
  \node[font=\scriptsize, thr, anchor=north west] at \BP{-1.95}{\bb-0.05} {$F_R$};
  \draw[thr, line width=1.2pt, -{Stealth[length=4pt,width=3pt]}]
       \BP{\lb}{0} -- \BP{\lb}{1.3};
  \node[font=\scriptsize, thr, anchor=south] at \BP{\lb+0.1}{0.75} {$F_B$};

  % 모멘트 팔 b — 선체 왼쪽 바깥
  \draw[ctrl, line width=0.45pt, {Stealth[length=2.6pt]}-{Stealth[length=2.6pt]}]
       \BP{-2.6}{0} -- \BP{-2.6}{\bb};
  \node[font=\scriptsize, ctrl, anchor=north] at \BP{-2.72}{0.5} {$b$};
  \draw[ctrl, line width=0.25pt, densely dotted] \BP{-2.7}{0}   -- \BP{-1.95}{0};
  \draw[ctrl, line width=0.25pt, densely dotted] \BP{-2.7}{\bb} -- \BP{-1.95}{\bb};

  % 모멘트 팔 l3 — 선체 아래 바깥
  \draw[ctrl, line width=0.45pt, {Stealth[length=2.6pt]}-{Stealth[length=2.6pt]}]
       \BP{0}{-1.75} -- \BP{\lb}{-1.75};
  \node[font=\scriptsize, ctrl, anchor=east] at \BP{0.37}{-1.85} {$l_3$};
  \draw[ctrl, line width=0.25pt, densely dotted] \BP{0}{-1.85}   -- \BP{0}{-1.12};
  \draw[ctrl, line width=0.25pt, densely dotted] \BP{\lb}{-1.85} -- \BP{\lb}{-1.12};

  % 양의 요 모멘트 — 선체 밖 오른쪽, 시계방향
  \draw[errc, line width=0.6pt, -{Stealth[length=3.4pt,width=2.4pt]}]
       (30.9,36.8) arc (50:-10:48mm);
  \node[font=\scriptsize, errc, anchor=west] at (50,22) {$N$ — 양의 요 모멘트};
  \node[font=\scriptsize, errc, anchor=west] at (50,16) {(선수가 우현으로)};

  \node[anchor=north west, font=\scriptsize, text width=96mm] at (-50,-70)
       {\textbf{한 추진기가 내는 세 몫} \\[3pt]
        $F_L,\ F_R$ 은 \textbf{선체 앞뒤}로 민다. 그래서 $X$ 에 그대로 들어가고
        $Y$ 에는 한 톨도 들어가지 않는다. 중심선에서 $b$ 만큼 떨어져 있으므로
        돌림 $N$ 에는 $+b\,F_L$ 과 $-b\,F_R$ 로 들어간다 \\[3pt]
        $F_B$ 는 \textbf{선체 좌우}로 민다. 그래서 $Y$ 에 그대로 들어가고 $X$ 에는 없다.
        무게중심에서 앞으로 $l_3$ 만큼 나가 있으므로 $N$ 에는 $+l_3 F_B$ 로 들어간다 \\[3pt]
        $\;\to\;$ 이 여섯 줄을 그대로 옮겨 적은 것이 배분행렬이다. 외울 것이 없다};
\end{scope}

% ===========================================================================
% (나) 옆으로만 밀고 싶을 때
% ===========================================================================
\begin{scope}[shift={(152,0)}]
  \node[font=\small\bfseries, anchor=north west] at (-50,66)
       {(나) 옆으로만 밀려면 — 선수추진기 혼자서는 안 된다};

  \hullplan

  % 선수추진기 — 옆으로 미는 힘
  \draw[thr, line width=1.6pt, -{Stealth[length=5pt,width=3.6pt]}]
       \BP{\lb}{0} -- \BP{\lb}{1.55};
  \node[font=\scriptsize, thr, anchor=south] at \BP{\lb+0.1}{0.9} {$F_B = +250$ N};

  % 좌우 추진기 — 서로 반대로
  \draw[thr, line width=1.2pt, -{Stealth[length=4pt,width=3pt]}]
       \BP{-1.8}{-\bb} -- \BP{-2.5}{-\bb};
  \node[font=\scriptsize, thr, anchor=north] at \BP{-2.6}{-\bb} {$F_L = -91.3$ N};
  \draw[thr, line width=1.2pt, -{Stealth[length=4pt,width=3pt]}]
       \BP{-1.8}{\bb} -- \BP{-1.1}{\bb};
  \node[font=\scriptsize, thr, anchor=north] at \BP{-2.0}{\bb} {$F_R = +91.3$ N};

  % 선수추진기가 내는 돌림 — 양(시계방향), 선체 밖 오른쪽
  \draw[errc, line width=0.6pt, -{Stealth[length=3.4pt,width=2.4pt]}]
       (30.9,36.8) arc (50:0:48mm);
  \node[font=\scriptsize, errc, anchor=west] at (50,26) {$+\,l_3 F_B = +187.5$ N$\cdot$m};

  % 좌우 추진기가 내는 돌림 — 음(반시계), 선체 밖 왼쪽
  \draw[plant, line width=0.6pt, -{Stealth[length=3.4pt,width=2.4pt]}]
       (-47.3,-8.3) arc (190:230:48mm);
  \node[font=\scriptsize, plant, anchor=east] at (-50,-20)
       {$b\,(F_L - F_R) = -187.5$ N$\cdot$m};

  \node[anchor=north west, font=\scriptsize, text width=96mm] at (-50,-70)
       {\textbf{왜 세 개가 모두 움직이는가} \\[3pt]
        배를 옆으로만 밀고 싶다 — 즉 $X = 0$, $Y = 250$ N, $N = 0$ 이다 \\[3pt]
        $Y$ 를 낼 수 있는 것은 선수추진기뿐이므로 $F_B = 250$ N 이 정해진다.
        그런데 선수추진기는 \textbf{앞쪽에} 있어, 혼자 밀면 배가 제자리에서 돌아 버린다
        ($+187.5$ N$\cdot$m) \\[3pt]
        그 돌림을 지우려고 좌우 추진기가 \textbf{서로 반대로} 민다.
        한쪽은 앞으로 $+91.3$ N, 다른 쪽은 뒤로 $-91.3$ N — 합이 $0$ 이라 $X$ 는 그대로 $0$ 이고,
        돌림만 $-187.5$ N$\cdot$m 로 나와 정확히 상쇄된다 \\[3pt]
        $\;\to\;$ 옆걸음 접안에서 \textbf{세 추진기가 모두 일한다}는 말은 이 뜻이다};
\end{scope}

% ===========================================================================
% 아래 띠 — 배분행렬과 그 역
% ===========================================================================
\draw[wire, soft] (-50,-118) -- (248,-118);

\node[anchor=north west, font=\scriptsize, text width=136mm] at (-50,-123)
     {\textbf{위 그림을 그대로 옮긴 것이 배분행렬이다} \\[4pt]
      $\boldsymbol\tau = T\,\mathbf{u}$, \qquad
      $\begin{bmatrix} X \\ Y \\ N \end{bmatrix}
       = \begin{bmatrix} 1 & 1 & 0 \\ 0 & 0 & 1 \\ b & -b & l_3 \end{bmatrix}
         \begin{bmatrix} F_L \\ F_R \\ F_B \end{bmatrix}$ \\[4pt]
      첫 줄 — 앞뒤로 미는 것은 좌우 추진기 둘뿐 \\[2pt]
      둘째 줄 — 옆으로 미는 것은 선수추진기 하나뿐 \\[2pt]
      셋째 줄 — 돌리는 몫은 각자의 \textbf{모멘트 팔}을 곱한 것};

\node[anchor=north west, font=\scriptsize, text width=136mm] at (112,-123)
     {\textbf{거꾸로 푸는 것이 추력 배분이다} \\[4pt]
      제어기가 낸 $X,\ Y,\ N$ 을 추진기 셋의 힘으로 되돌린다 —
      $\mathbf{u} = T^{-1}\boldsymbol\tau$. 손으로 풀면 세 줄이다 \\[4pt]
      $F_B = Y$ \qquad (둘째 줄에서 바로 나온다) \\[2pt]
      $M = N - l_3 F_B$ \qquad (선수추진기가 이미 쓴 돌림을 빼고 남은 몫) \\[2pt]
      $F_L = \dfrac{X}{2} + \dfrac{M}{2b}$, \qquad
      $F_R = \dfrac{X}{2} - \dfrac{M}{2b}$ \\[4pt]
      앞으로 미는 몫은 둘이 나눠 지고, 돌리는 몫은 둘이 \textbf{반대로} 진다};

\end{tikzpicture}
\end{document}
